% chktex-file 1
% chktex-file 44
% chktex-file 45
% chktex-file 36
% chktex-file 3
% chktex-file 8
% chktex-file 18
% chktex-file 12
% chktex-file 13

\documentclass[twoside]{article}
\usepackage{graphicx,fancyhdr,amsmath,amssymb,amsthm,subfig,url,hyperref}
\usepackage[margin=1in]{geometry}
\newtheorem{theorem}{Theorem}


%---------------- IMPORTANT: Student and Homework Information -------------

%%% PLEASE FILL THIS OUT WITH YOUR INFORMATION
\newcommand{\myname}{Your Name: Muhammad Mujtaba}
\newcommand{\myid}{28100480}
\newcommand{\hwNo}{Assignment 7}
%%% END



\fancypagestyle{plain}{}
\pagestyle{fancy}
\fancyhf{}
\fancyhead[RO,LE]{\sffamily\bfseries\large LUMS}
\fancyhead[LO,RE]{\sffamily\bfseries\large CS 210 Discrete Mathematics}
\fancyfoot[LO,RE]{\sffamily\bfseries\large \myname: \myid @lums.edu.pk}
\fancyfoot[RO,LE]{\sffamily\bfseries\thepage}
\renewcommand{\headrulewidth}{1pt}
\renewcommand{\footrulewidth}{1pt}

%--------------------- This is the title of the document. DO NOT CHANGE IT ------------------------

\title{CS 210 \hwNo}
\author{\myname \qquad Student ID:\myid}
\date{}

%--------------------------------- AFTER Entering the Student and Homework Information, write your answers below  ----------------------------------

\begin{document}
\maketitle

\section*{Question 1}
\begin{enumerate}
    \item[(a)] Permutations of BANANAS:\[\frac{7!}{3!\,2!}=420.\]
    \item[(b)] Find the 244th word in lexicographic order:
    \begin{itemize}
        \item 1. Words starting with A: The remaining letters are \{B, N, A, N, A, S\}. \[ \text{Count} = \frac{6!}{2!2!1!1!} = 180 \]
        \item 2. Words starting with B: The remaining letters are \{A, N, A, N, A, S\}. \[ \text{Count} = \frac{6!}{3!2!1!} = 60 \]
        \item 3. Determine the target range: \\
        Cumulative total = $180 + 60 = 240$.
        We need the 244th word. Since $244 > 240$, the word starts with the next letter, N. \\
        We are looking for the 4th word starting with N ($244 - 240 = 4$).
        \item 4. Finding the 4th word starting with N:
        The remaining letters are \{A, A, A, B, N, S\}.
        \begin{itemize}
            \item Try prefix NA: Remaining letters \{A, A, B, N, S\}. 
            Permutations: $\frac{5!}{2!} = 60$. 
            Since $4 \le 60$, the word starts with NA.
            \item Try prefix NAA: Remaining letters \{A, B, N, S\}. 
            Permutations: $4! = 24$.
            Since $4 \le 24$, the word starts with NAA.
            \item Try prefix NAAA: Remaining letters \{B, N, S\}. 
            Permutations: $3! = 6$.
            Since $4 \le 6$, the word starts with NAAA.
        \end{itemize}
        
        \item 5. Final Step:
        We need the 4th permutation of the remaining letters \{B, N, S\}.
        Sorted lexicographically: B, N, S.
        \begin{enumerate}
            \item BNS
            \item BSN
            \item NBS
            \item NSB
        \end{enumerate}
        Appending this suffix to the prefix NAAA gives the final word.
    \end{itemize}
    
    \textbf{Answer:} The 244th word is \textbf{NAAANSB}.
\end{enumerate}

\newpage

\section*{Question 2}
\begin{enumerate}
    \item[(a)] A five-digit palindrome has form \(abcba\).  \[9\cdot 10\cdot 10 = 900.\]
    \item[(b)] Odd palindromes: last digit must be odd → 5 choices: \[5\cdot 10 \cdot 10 = 500.\] 
               Even palindromes:\[4\cdot 10 \cdot 10 = 400.\]
\end{enumerate}

\section*{Question 3}
\begin{enumerate}
    \item[(a)]  Coefficient of \(x^5\) in \((1+2x)^{11}\) \[\binom{11}{5}2^5 = 462\cdot 32 = 14784.\]
    \item[(b)]  Coefficient of \(x^7\) in \((2-x)^{14}\) \[-\binom{14}{7}2^7 = -3432\cdot 128 = -439296.\]
    \item[(c)]  Coefficient of \(x^9 y^9\) in \((4x-3y)^{18}\) \[\binom{18}{9}(4^9)(-3)^9.\]
\end{enumerate}

\section*{Question 4}
We count binary string of length \(n\) with one designated 1.
Counting by number of ones: \[\sum_{k=1}^{n} k\binom{n}{k}.\]
Counting by position of designated 1: \[n\cdot 2^{n-1}.\]
Thus: \[\sum_{k=1}^{n} k\binom{n}{k} = n2^{n-1}.\]

\newpage

\section*{Question 5}
Prove the hexagon identity: \[ \binom{n-1}{k-1}\binom{n}{k+1}\binom{n+1}{k} = \binom{n-1}{k}\binom{n}{k-1}\binom{n+1}{k+1} \]
Proof: We expand the binomial coefficients using \[\binom{n}{r} = \frac{n!}{r!(n-r)!}\]
LHS Expansion: \[\frac{(n-1)!}{(k-1)!(n-k)!} \cdot \frac{n!}{(k+1)!(n-k-1)!} \cdot \frac{(n+1)!}{k!(n+1-k)!} \]
Grouping numerators and denominators:\[ \text{LHS} = \frac{(n-1)! n! (n+1)!}{ [(k-1)! (k+1)! k!] \cdot [(n-k)! (n-k-1)! (n-k+1)!] }\]

RHS Expansion: \[\frac{(n-1)!}{k!(n-1-k)!} \cdot \frac{n!}{(k-1)!(n-k+1)!} \cdot \frac{(n+1)!}{(k+1)!(n-k)!}\]
Grouping numerators and denominators: \[ \text{RHS} = \frac{(n-1)! n! (n+1)!}{ [k! (k-1)! (k+1)!] \cdot [(n-k-1)! (n-k+1)! (n-k)!] }\]
Conclusion:
The numerators are identical. The sets of factorials in the denominators are just rearranged but contain the exact same terms. Therefore, LHS = RHS.

\section*{Question 6}
Suppose 6 distinct devices (3 distinct interface cards, printer, mouse, modem) must be plugged into 8 distinct slots.
\begin{enumerate}
    \item[(a)] Total number of ways \\
               We choose an ordered set of 6 slots from 8:\[ {}_{8}P_{6} = 8\cdot7\cdot6\cdot5\cdot4\cdot3 = 20160. \]
    \item[(b)] No two interface cards adjacent \\
               We choose an ordered set of 6 slots from 8:\[ {}_{8}P_{6} = 8\cdot7\cdot6\cdot5\cdot4\cdot3 = 20160. \]
               Number of ways to choose 3 non-adjacent positions among 8 is: \[ \binom{8-3+1}{3} = \binom{6}{3} = 20. \]
               Arrange the 3 cards: \[3! = 6. \]
               Choose 3 of the remaining 5 slots for the 3 remaining devices: \[ {}_{5}P_{3} = 5\cdot4\cdot3 = 60. \]
               Thus total: \[ 20 \cdot 6 \cdot 60 = 7200. \]
\end{enumerate}

\section*{Question 7}
Since the teams are indistinguishable (unlabeled), we can categorize the valid distributions based on the integer partitions of 7 into 3 parts (i.e., the possible sizes of the teams). There are four distinct cases:

\begin{enumerate}
    \item Case 1: Team sizes are \{5, 1, 1\} \\
    We first select 5 people for the large team. The remaining 2 people automatically form two single-person teams. Since these two teams are identical in size (1) and unlabelled, their order does not matter. \[ \text{Ways} = \binom{7}{5} = 21 \]
    \item Case 2: Team sizes are \{4, 2, 1\} \\
    Here, all team sizes are distinct (4, 2, and 1), which makes the groups distinguishable by size.
    \begin{itemize}
        \item Choose 4 people for the first team: $\binom{7}{4} = 35$
        \item Choose 2 people from the remaining 3: $\binom{3}{2} = 3$
        \item The last person takes the final spot: $\binom{1}{1} = 1$
    \end{itemize}
    \[ \text{Ways} = 35 \times 3 \times 1 = 105 \]

    \item Case 3: Team sizes are \{3, 2, 2\} \\
    First, we choose 3 people for the size-3 team: $\binom{7}{3} = 35$.
    The remaining 4 people must be partitioned into two unlabelled sets of size 2. We choose 2 people for one group, but since the two groups are indistinguishable, we divide by \[2!\] to correct for overcounting.
    \[ \text{Ways} = \binom{7}{3} \times \frac{\binom{4}{2}}{2!} = 35 \times \frac{6}{2} = 35 \times 3 = 105 \]

    \item Case 4: Team sizes are \{3, 3, 1\} \\
    First, we choose 1 person for the size-1 team: $\binom{7}{1} = 7$.
    The remaining 6 people must be partitioned into two unlabelled sets of size 3. Similar to the previous case, we divide by \[2!.\]
    \[ \text{Ways} = \binom{7}{1} \times \frac{\binom{6}{3}}{2!} = 7 \times \frac{20}{2} = 7 \times 10 = 70 \]
\end{enumerate}

\textbf{Total Ways:}
Summing the results from the four cases:
\[ 21 + 105 + 105 + 70 = 301 \]

\textbf{Answer:} There are \textbf{301} ways to choose the teams.
\end{document}